\section{Introduction}\label{sec:introduction}

Basic sanitation is essential infrastructure for public health, environmental protection, and socioeconomic development \cite{barros1995}.
Nevertheless, the deployment of sanitary sewer networks remains uneven in Brazil, especially in small municipalities, peripheral areas, and regions with limited technical and budgetary capacity.
In these contexts, the initial planning phase often becomes a bottleneck: before construction, it is necessary to collect data, interpret terrain, define preliminary routes, estimate inspection structures, and compare implementation alternatives.

Traditional preliminary design depends on specialized knowledge, engineering software, and intensive manual intervention.
This combination makes the process expensive and slow, hindering feasibility studies and delaying the expansion of infrastructure in vulnerable areas.
Prior work on sewer network optimization, urban management, and computational modeling indicates that graph techniques, search algorithms, and integration with geospatial data can reduce this initial effort when embedded in tools that can be used by domain specialists \cite{weng2014,haghighi2012,melo2019,rodrigues2020}.

\method was conceived in this context as a platform to automate and support steps of preliminary sanitary sewer network planning.
The research builds on a method previously applied in a case study and evolves it into a functional web application, with a more robust architecture, an interactive editor, and an algorithmic pipeline capable of processing street networks enriched with elevation data.
The proposal does not replace executive engineering design; instead, it offers a computational basis to accelerate early conception and analysis.

This paper describes the current development of \method as a computational contribution applied to preliminary sanitary sewer network planning.
The remaining sections present the theoretical background, related systems, proposed method, architecture and implementation, current validation status, complementary experiments, and limitations of a version that is still undergoing systematic experimental validation.
